\section{外代数的函子性}

记号约定:$A$是交换环.

\begin{proposition}{外代数的自然性}{naturality-of-exterior-algebra}
	设$M,N$是$A$-模,$u\in\Hom_{A}\left(M,N\right)$,则存在唯一的$u^{\prime}\in\Hom_{A-\mathsf{Alg}}\left(\bigwedge\left(M\right),\bigwedge\left(N\right)\right)$使下图交换.
	\begin{center}
		\begin{tikzcd}
			M & N \\
			{\bigwedge\left(M\right)} & {\bigwedge\left(N\right)}
			\arrow["u", from=1-1, to=1-2]
			\arrow["{\phi_{M}^{\prime\prime}}"', from=1-1, to=2-1]
			\arrow["{\phi_{N}^{\prime\prime}}", from=1-2, to=2-2]
			\arrow["{\exists!u^{\prime\prime}}"', dashed, from=2-1, to=2-2]
		\end{tikzcd}
	\end{center}
\end{proposition}

\begin{definition}{映射的典范扩张}{}
	沿用命题(\cref{prop:naturality-of-exterior-algebra})的设定.我们把$u^{\prime\prime}$记作$\bigwedge\left(u\right)$,称之为$u$到$\bigwedge\left(M\right)$的典范扩张.
\end{definition}

\begin{proposition}{外代数的函子性}{}
	设$M,N$是$A$-模,$u\in\Hom_{A}\left(M,N\right),v\in\Hom_{A}\left(N,P\right)$,则下列图表交换.
	\begin{center}
		\begin{tikzcd}
			M & N & P \\
			{\bigwedge\left(M\right)} & {\bigwedge\left(N\right)} & {\bigwedge\left(P\right)}
			\arrow["u", from=1-1, to=1-2]
			\arrow["{\bigwedge(v\circ u)}", curve={height=-18pt}, from=1-1, to=1-3]
			\arrow["{\phi_{M}^{\prime}}"', from=1-1, to=2-1]
			\arrow["v", from=1-2, to=1-3]
			\arrow["{\phi_{N}^{\prime\prime}}", from=1-2, to=2-2]
			\arrow["{\phi_{P}^{\prime\prime}}", from=1-3, to=2-3]
			\arrow["{\bigwedge\left(u\right)}"', from=2-1, to=2-2]
			\arrow["{\bigwedge\left(v\right)\circ\bigwedge\left(u\right)}"', curve={height=24pt}, from=2-1, to=2-3]
			\arrow["{\bigwedge\left(v\right)}"', from=2-2, to=2-3]
		\end{tikzcd}
	\end{center}
\end{proposition}

\begin{remark}{}{}
	设$n\in\N$.记$\bigwedge\limits^{n}\left(u\right):\bigwedge\limits^{n}\left(M\right)\to\bigwedge\limits^{n}\left(N\right),x\mapsto\bigwedge\left(u\right)\left(x\right)$,那么对$M$中的任一有限序列$\left(x_{i}\right)_{i=1}^{n}$有
	\begin{align*}
		\bigwedge\limits^{n}\left(u\right)\left(x_{1}\wedge\ldots\wedge x_{n}\right)=u\left(x_{1}\right)\wedge\ldots\wedge u\left(x_{n}\right).
	\end{align*}
	此外,$\bigwedge\limits^{0}\left(u\right)$等同于$\id_{A}$.
\end{remark}

\begin{proposition}{外代数上的典范扩张的性质}{}
	设$M,N$是$A$-模,$u\in\Hom_{A}\left(M,N\right)$.若$u$是满射,则
	\begin{enumerate}
		\item $\bigwedge\left(u\right)$是满射.
		\item $\ker\left(\bigwedge\left(u\right)\right)$是$\ker\left(u\right)$在$\bigwedge\left(M\right)$中生成的双边理想.
	\end{enumerate}
\end{proposition}

\begin{remark}{}{}
	如果$u$是单射,那么$\bigwedge\left(u\right)$一般不是单射.
\end{remark}

\begin{proposition}{}{}
	设$M$是$A$-模,$N,P$是$M$的子模,$N+P$是$M$的直因子,$N\cap P$是$N$和$P$各自的直因子.考虑典范单射
	\begin{align*}
		i:N\to M,j:P\to M,k:N\cap P\to M.
	\end{align*}
	\begin{enumerate}
		\item $\bigwedge\left(i\right),\bigwedge\left(j\right),\bigwedge\left(k\right)$都是单射.
		\item $\im\left(\bigwedge\left(k\right)\right)=\im\left(\bigwedge\left(i\right)\right)\cap\im\left(\bigwedge\left(j\right)\right)$.
	\end{enumerate}
\end{proposition}

\begin{remark}{}{}
	如果$A$是域,那么``$N\cap P$是$N$和$P$各自的直因子''自动成立.
\end{remark}

\begin{corollary}{}{}
	设$K$是域,$M$是$K$-向量空间,则对每个$z\in\bigwedge\left(M\right)$,$M$都有一个最小的有限维子空间$N$满足$z\in\bigwedge\left(N\right)$.
\end{corollary}

\begin{definition}{与$z$关联的向量子空间}{}
	设$K$是域,$M$是$K$-向量空间,$z\in\bigwedge\left(M\right)$.$M$的满足$z\in\bigwedge\left(N\right)$的最小有限维子空间$N$称为$M$中与$z$关联的向量子空间.
\end{definition}